% Unnummeriert, damit die sieben Teile des Musters die Nummern 1 bis 7 tragen.
\section*{Rahmen und Methode}
\addcontentsline{toc}{section}{Rahmen und Methode}
\markright{Rahmen und Methode}
\label{sec:methode}

\subsection*{Gegenstand}

GEA Group Aktiengesellschaft, D\"usseldorf, ist ein Maschinen- und Anlagenbauer f\"ur die
Nahrungsmittel-, Getr\"anke- und Pharmaindustrie. Der Bericht betrachtet die
Gesch\"aftsjahre 2019 bis 2025 sowie das erste Halbjahr 2026 und entscheidet \"uber den
Kauf der Aktie zum Kurs von \je{\Kurs}{EUR} (\KursDatum).

\subsection*{Quellen}

Prim\"arquellen sind die Gesch\"aftsberichte 2020 bis 2025 und der Halbjahresfinanzbericht
2026, jeweils in der englischen Originalfassung und lokal abgelegt. Jede Seitenangabe
bezeichnet die \emph{gedruckte} Seite, nicht die Seitenzahl der PDF-Datei; beide
unterscheiden sich bei GEA um mehrere Seiten. Sekund\"arquellen sind als solche
gekennzeichnet: Kurse (Yahoo Finance), der Analystenkonsens (MarketScreener), die
Directors' Dealings (EQS-Meldungen nach Art.~19 MAR) und die H\"urde (MSCI).

\subsection*{Wie die Zahlen in dieses Dokument kommen}

Die Gesch\"aftsberichte werden seitenweise zerlegt; ein lokales Sprachmodell benennt
lediglich die \emph{Zeilenetiketten} der Abschl\"usse, die Werte liest Programmcode aus
dem Textlayer derselben Seite. Danach pr\"uft \texttt{scripts/pruefe\_seiten.py} jeden
gespeicherten Wert gegen den Text der Seite, die er als Beleg nennt:
\GeprueftAnzahl{} von \GeprueftAnzahl{} Werten sind so belegt. Der Flie{\ss}text enth\"alt
keine Zahl von Hand; \texttt{scripts/check\_literals.py} bricht den Bau ab, sobald doch
eine erscheint.

Diese Pr\"ufung erkennt eine falsch zugeordnete Zahl nicht, wenn die Zahl auf der Seite
steht. Gegen diesen Fall stehen zwei weitere Pr\"ufungen: die Kette \"uber zwei Berichte
(jeder Bericht druckt das Vorjahr, beide m\"ussen \"ubereinstimmen) und der Vergleich der
Kennzahlenseite mit der Kapitalflussrechnung.

\subsection*{Die drei Gr\"o{\ss}en, auf denen der Bericht ruht}

\textbf{Freier Zufluss.} Operativer Mittelzufluss der fortgef\"uhrten T\"atigkeit,
vermindert um Investitionen und um die Tilgung von Leasingverbindlichkeiten. GEA weist
einen eigenen \enquote{Free Cash Flow} aus, der 2025 mit \Mio{\FreierCashflowGea}{EUR}
\"uber dieser Gr\"o{\ss}e liegt\Q{GeaGB}{2}; der Bericht rechnet mit der eigenen Gr\"o{\ss}e,
weil nur sie den Leasingabfluss enth\"alt, den ein Eigent\"umer ebenso tragen muss wie
eine Investition.

\textbf{H\"urde.} \Pz{\Huerde} p.\,a.\ Das ist die Rendite des MSCI World (Gross Returns,
USD) \"uber zehn Jahre nach dem Factsheet zum 31.08.2026\QW{MSCI World Index, Index
Factsheet 31.08.2026}{quelle:msci}{19.09.2026}. Eine Anlage in GEA muss diese Rendite
schlagen, sonst ist der Indexfonds die bessere Wahl. Als Sensitivit\"at dient die
Rendite seit Auflage des Index (1987) von \Pz{\HuerdeLang}; sie f\"uhrt zu einem
h\"oheren Schwellenkurs und ist damit die f\"ur GEA g\"unstigere Annahme.

\textbf{Endwert.} Der Buchwert des Eigenkapitals zum 31.12.2025,
\Mio{\Eigenkapital}{EUR}\Q{GeaGB}{312}, entspricht \je{\BuchwertJeAktie}{EUR} je Aktie.

\annahme{Alle Rechnungen sind Vorsteuerrechnungen auf Ebene des Unternehmens. Der
Zufluss ist der Zufluss an das Unternehmen, nicht der nach Abgeltungsteuer beim Anleger
verbleibende Betrag. Wer die Anlage in einem steuerpflichtigen Depot h\"alt, muss die
ausgewiesenen Renditen entsprechend k\"urzen; die H\"urde (Indexrendite) w\"are dann
ebenfalls zu k\"urzen, sodass der Vergleich der Optionen untereinander bestehen bleibt.}

\annahme{Der freie Zufluss wird in allen Szenarien \emph{konstant} fortgeschrieben. Das
ist keine Prognose, sondern der Verzicht auf eine. Der Bericht holt den fehlenden
Wachstumsterm in Abschnitt~\ref{sec:votum} als Umkehrung zur\"uck: statt Wachstum zu
sch\"atzen, rechnet er aus, welches Wachstum der heutige Kurs bereits voraussetzt, und
h\"alt es gegen die eigene Reihe und gegen die Ziele des Unternehmens.}
